% Figure: an I-section decomposed into three rectangles, with the
% centroidal axis and the parallel-axis offsets d_i drawn.
\begin{figure}[htbp]
\centering
\begin{tikzpicture}[>=latex, x=1.0cm, y=1.0cm]
  % I-section: top flange, web, bottom flange (symmetric -> centroid at mid-height)
  % flange 200x20 -> draw width 4, thickness 0.4; web 200x10 -> width 1, height 4
  % overall depth 240 -> 4.8
  \def\fw{4.0} \def\ft{0.4} \def\wh{4.0} \def\ww{1.0}
  % bottom flange
  \fill[engblue!18, draw=engblue, thick] (-\fw/2,0) rectangle (\fw/2,\ft);
  % web
  \fill[engblue!10, draw=engblue, thick] (-\ww/2,\ft) rectangle (\ww/2,\ft+\wh);
  % top flange
  \fill[engblue!18, draw=engblue, thick] (-\fw/2,\ft+\wh) rectangle (\fw/2,2*\ft+\wh);
  % centroidal axis (mid-height by symmetry)
  \pgfmathsetmacro\yc{\ft+\wh/2}
  \draw[engred, very thick, dashed] (-\fw/2-0.6,\yc) -- (\fw/2+0.6,\yc);
  \node[engred, anchor=west, font=\scriptsize] at (\fw/2+0.6,\yc) {centroidal axis};
  % offset for top flange centroid
  \pgfmathsetmacro\ytf{\ft+\wh+\ft/2}
  \pgfmathsetmacro\ymid{(\yc+\ytf)/2}
  \pgfmathsetmacro\xlab{\fw/2+0.3}
  \pgfmathsetmacro\xarr{\fw/2+0.25}
  \draw[engamber, thin] (-\fw/2,\ytf) -- (\fw/2,\ytf);
  \draw[<->, engamber, thin] (\xarr,\yc) -- (\xarr,\ytf);
  \node[engamber, anchor=west, font=\scriptsize] at (\xlab,\ymid) {$d$};
  % centroids of each rectangle
  \fill[engblack] (0,\ft/2) circle (1.5pt);
  \fill[engblack] (0,\yc) circle (1.5pt);
  \fill[engblack] (0,\ytf) circle (1.5pt);
  % dimension labels
  \node[font=\scriptsize, anchor=north] at (0,-0.1) {$b_f$};
  \draw[<->, enggray, thin] (-\fw/2,-0.05) -- (\fw/2,-0.05);
  \node[font=\scriptsize, anchor=east] at (-\fw/2-0.15,\yc) {$h$};
  \draw[<->, enggray, thin] (-\fw/2-0.1,0) -- (-\fw/2-0.1,2*\ft+\wh);
\end{tikzpicture}
\caption[Composite I-section and the parallel-axis theorem]{An I-section
treated as three rectangles for the second-moment calculation. Each
rectangle contributes its own centroidal second moment $bh^3/12$ plus a
transfer term $A_id_i^2$, where $d_i$ is the distance from the
rectangle's own centroid to the composite centroidal axis. For the
symmetric section the composite centroid sits at mid-height, so the web's
transfer term vanishes ($d=0$) and the two flanges contribute the bulk of
the stiffness through their large $Ad^2$ terms. This is the quantitative
reason a designer pushes material to the extremes of the section: the
transfer term grows with the square of the distance, so a kilogram of
steel in the flange buys far more bending stiffness than the same
kilogram in the web. The parallel-axis theorem turns a hard integral into
a sum of rectangle areas and offsets done by hand.}
\label{fig:vol03ch02:centroid-composite}
\end{figure}
